% generated by bibm/supplement/make_sm_tables.py -- do not edit; sources: shuffled_neg.csv, alt_model_shuffle_hashseed.csv
\begin{table}[!htbp]
\caption{Dinucleotide-shuffled negatives: the paper's draw and the range over 16 \task{PYTHONHASHSEED} draws of the shuffle. Source: \src{shuffled\_neg.csv},\enspace \src{alt\_model\_shuffle\_hashseed.csv}.}
\label{tab:shuffle}
\centering
\scriptsize
\setlength{\tabcolsep}{3pt}
\begin{tabular}{lllllll}
\toprule
 & \multicolumn{2}{c}{TATA motif-only AUROC} &  &  &  &  \\
Data & real neg. & shuffled (range) & GC-content AUROC & $k$-mer AUROC (range) & $k$-mer MCC (range) & = paper \\
\midrule
\task{cohn} & 0.371 & 0.551 (0.549--0.554) & 0.500 & 0.955 (0.951--0.955) & 0.769 (0.759--0.776) & 0/16 \\
\task{nt\_enhancers} & 0.393 & 0.507 (0.503--0.510) & 0.500 & 0.856 (0.839--0.879) & 0.545 (0.520--0.627) & 0/16 \\
control & 0.539 & 0.749 (0.746--0.754) & 0.500 & 0.989 (0.988--0.996) & 0.922 (0.906--0.937) & 0/16 \\
\bottomrule
\end{tabular}
\par\smallskip
\begin{minipage}{\linewidth}\footnotesize\raggedright
Shuffles of the positives remove the mono- and dinucleotide composition gap by construction (GC-content AUROC 0.500 on every task; the GC-content AUROC is that of per-sequence GC content, not a motif score). The negative-class TATA skew disappears (\task{cohn} TATA AUROC slightly above 0.5 in every draw), while the positives remain separable from their own shuffles by higher-order structure; for the control, whose TATA is a positive-class motif, the TATA AUROC rises. All sequences are scanned. The shuffle depends on \task{PYTHONHASHSEED}; the paper's exact draw was not recovered by the 16-seed scan.
\end{minipage}
\end{table}
